\documentclass[12pt]{article}
\usepackage[utf8]{inputenc}
\usepackage[T1]{fontenc}
\usepackage{hyperref}
\title{Ngardi Language Resources for K-12 Curriculum Development}
\author{Compiled by Claude Sonnet 5.5 (claude-sonnet-5-5) for Andr\'as Kornai}
\date{}
\begin{document}
\maketitle

\section{Introduction}
Ngardi is a Pama-Nyungan language of the western Tanami Desert and the north of the Great Sandy Desert. Glottolog calls it extinct. Australian Geographic reports that only a handful of people speak it fluently and that children no longer learn it as a first language. Ngardi people live in the Kutjungka communities of Mulan and Billiluna, in Halls Creek, Ringers Soak and Kununurra, and in Yuendumu (Northern Territory), according to NIT.

\section{Primary Educational Resources}
No school curriculum, primer or classroom programme specific to Ngardi was found in the search results.

\section{Linguistic and Reference Materials}
\begin{itemize}
\item \emph{Ngardi to English Dictionary} (Aboriginal Studies Press, 2025), compiled by Tom Ennever, Marie Mudgedell, Tjama Napanangka and Lee Cataldi. It has an alphabet, an English finder list, a verb list, a grammar guide, kinship information, example sentences and a history of the language and Country. Listings: \url{https://readings.com.au/product/9781922102850} and \url{https://amplifybookstore.com/products/ngardi-to-english-dictionary}
\item Publisher resource sheet: \url{https://newsouthbooks.com.au/assets/media/resources/feb25a_rs.pdf}
\item AIATSIS AUSTLANG entry: \url{https://collection.aiatsis.gov.au/austlang/language/c2}. It lists word lists by Capell (1952) and Wafer (1980) and records a Warlpiri-based practical orthography.
\item South Australian Museum archive, Ngardi (NT): \url{https://www.samuseum.sa.gov.au/collection/archives/language_groups/ngardi}
\end{itemize}

\section{Religious and Community Resources}
\begin{itemize}
\item NIT article on the dictionary, ``Spoken since time immemorial'': \url{https://nit.com.au/13-05-2025/17915/spoken-since-time-immemorial-new-dictionary-preserves-ngardi-language}
\item Australian Geographic, ``Deadly Dictionaries: Ngardi'': \url{https://australiangeographic.com.au/history-culture/2026/08/deadly-dictionaries-ngardi/}
\end{itemize}

\section{Digital and Audio Resources}
Territory Stories (NT Library) holds Living Archive of Aboriginal Languages material for the neighbouring Warlpiri area, for example \emph{Miyi Lajamanu-wana}: \url{https://territorystories.nt.gov.au/10070/881470}. It is not specific to Ngardi. A second Territory Stories item also appeared in the results: \url{https://territorystories.nt.gov.au/10070/891124/0/2}. Its content was not verified. No Ngardi audio or software was found.

\section{Curriculum Development Context}
The 2025 dictionary is the main modern resource. Its grammar guide, example sentences and alphabet could support a heritage and revitalization curriculum. Any such curriculum should be developed with Ngardi elders and with the communities named above. Warlpiri-area bilingual materials may serve as models, but they are not Ngardi.

\end{document}